\documentclass[../../../include/open-logic-section]{subfiles}

\begin{document}
\olfileid{sth}{card-arithmetic}{expotough}

\olsection[Sawetara Panyederhanaan]{Sawetara Panyederhanaan Pemangkatan Kardinal}

Sanadyan netepaké $\canonord$ rada njlimet, asilé
pratelan kang migunani babagan panambahan lan ping-pingan
kardinal, \olref[simp]{cardplustimesmax}. Nanging pemangkatan
transfinit ora bisa disederhanakaké kanthi cara kang padha
gampangé. Kanggo nerangaké sebabé, kita miwiti kanthi asil
kang ngluwihi pola kang wis dikenal ing kasus winates
(sanadyan buktiné ana ing tingkat abstraksi kang dhuwur):

\begin{prop}\ollabel{simplecardexpo}
$\cardexpo{\cardfont{a}}{\cardfont{b} \cardplus \cardfont{c}} =
\cardexpo{\cardfont{a}}{\cardfont{b}} \cardtimes
\cardexpo{\cardfont{a}}{\cardfont{c}}$ lan
$\cardexpo{(\cardexpo{\cardfont{a}}{\cardfont{b}})}{\cardfont{c}} =
\cardexpo{\cardfont{a}}{\cardfont{b} \cardtimes \cardfont{c}}$ kanggo
sembarang kardinal $\cardfont{a}, \cardfont{b}, \cardfont{c}$.
\end{prop}

\begin{proof}
Kanggo pratelan kapisan, gatekna fungsi $f \colon
(\cardfont{b}\disjointsum\cardfont{c}) \to \cardfont{a}$.
Saiki “pisahna” fungsi iki kanthi netepaké
$f_\cardfont{b}(\beta) = f(\beta, 0)$ kanggo saben
$\beta \in \cardfont{b}$, lan $f_\cardfont{c}(\gamma) = f(\gamma, 1)$
kanggo saben $\gamma \in \cardfont{c}$. Pemetaan $f \mapsto
\tuple{f_{\cardfont{b}}, f_\cardfont{c}}$ iku !!a{bijection}
$\funfromto{\cardfont{b} \disjointsum \cardfont{c}}{\cardfont{a}} \to
(\funfromto{\cardfont{b}}{\cardfont{a}} \times
\funfromto{\cardfont{c}}{\cardfont{a}})$.

Kanggo pratelan kapindho, gatekna fungsi $f \colon \cardfont{c} \to
(\funfromto{\cardfont{b}}{\cardfont{a}})$; dadi kanggo saben $\gamma \in
\cardfont{c}$ kita nduwèni fungsi $f(\gamma) \colon \cardfont{b} \to
\cardfont{a}$. Saiki tetepna $f^*(\beta, \gamma) = (f(\gamma))(\beta)$
kanggo saben $\tuple{\beta, \gamma} \in \cardfont{b} \times \cardfont{c}$.
Pemetaan $f \mapsto f^*$ iku !!a{bijection}
$\funfromto{\cardfont{c}}{(\funfromto{\cardfont{b}}{\cardfont{a}})}
\to \funfromto{\cardfont{b} \times \cardfont{c}}{\cardfont{a}}$.
\end{proof}

Saiki kang \emph{kita karepake} yaiku cara gampang
kanggo ngitung $\cardexpo{\cardfont{a}}{\cardfont{b}}$
yèn kardinal-kardinalé tanpa wates. Iki sawijining
langkah kang becik menyang kana:

\begin{prop}\ollabel{cardexpo2reduct}
Yèn $2 \leq \cardfont{a} \leq \cardfont{b}$ lan $\cardfont{b}$
tanpa wates, mula $\cardexpo{\cardfont{a}}{\cardfont{b}} =
\cardexpo{2}{\cardfont{b}}$
\end{prop}

\begin{proof}
\begin{align*}
	\cardexpo{2}{\cardfont{b}} &\leq
	\cardexpo{\cardfont{a}}{\cardfont{b}}\text{, amarga $2 \leq \cardfont{a}$}\\
	&\leq \cardexpo{(2^\cardfont{a})}{\cardfont{b}}
	\text{, miturut \olref[opps]{cantorcor}}\\
	&= \cardexpo{2}{\cardfont{a} \cardtimes \cardfont{b}}
	\text{, miturut \olref{simplecardexpo}} \\
	&= \cardexpo{2}{\cardfont{b}}
	\text{, miturut \olref[simp]{cardplustimesmax}}
\end{align*}
\end{proof}

Ora pantes kita ngarep-arep panyederhanaan luwih adoh,
amarga $\cardfont{b} < \cardexpo{2}{\cardfont{b}}$ miturut
\olref[card-arithmetic][opps]{cantorcor}.
Nanging iki durung nerangaké apa kang bisa dikandhakake
babagan $\cardexpo{\cardfont{a}}{\cardfont{b}}$ nalika $\cardfont{b} <
\cardfont{a}$. Mesthiné, yèn $\cardfont{b}$ \emph{winates},
kita ngerti carané.

\begin{prop}
Yèn $\cardfont{a}$ tanpa wates lan $0<n \in \omega$, mula
$\cardexpo{\cardfont{a}}{n} = \cardfont{a}$
\end{prop}

\begin{proof}
$\cardexpo{\cardfont{a}}{n} = \cardfont{a} \cardtimes \cardfont{a}
\cardtimes \ldots \cardtimes \cardfont{a} = \cardfont{a}$, miturut \olref[simp]{cardplustimesmax}.
\end{proof}
\noindent
Kajaba iku, ing sawatara kasus liya, kita bisa matesi ukuran
$\cardexpo{\cardfont{a}}{\cardfont{b}}$:

\begin{prop}
Yèn $2 \leq \cardfont{b} < \cardfont{a} \leq
\cardexpo{2}{\cardfont{b}}$ lan $\cardfont{b}$ tanpa wates, mula
$\cardexpo{\cardfont{a}}{\cardfont{b}} = \cardexpo{2}{\cardfont{b}}$
\end{prop}

\begin{proof}
$\cardexpo{2}{\cardfont{b}}\leq \cardexpo{\cardfont{a}}{\cardfont{b}}
\leq \cardexpo{(\cardexpo{2}{\cardfont{b}})}{\cardfont{b}} =
\cardexpo{2}{\cardfont{b}\cardtimes\cardfont{b}} =
\cardexpo{2}{\cardfont{b}}$, kanthi panalaran kaya ing \olref{cardexpo2reduct}.
\end{proof}
\noindent
Nanging sawisé iki, prakarané dadi luwih rumit.

\end{document}
